\documentclass{article}
\usepackage[utf8]{inputenc}
\usepackage{hyperref}
\usepackage{listings}
\usepackage{xcolor}
\usepackage[margin=1in]{geometry}
\usepackage{tikz}
\usetikzlibrary{shapes.geometric, arrows.meta, positioning}

\definecolor{codegreen}{rgb}{0,0.6,0}
\definecolor{codegray}{rgb}{0.5,0.5,0.5}
\definecolor{codepurple}{rgb}{0.58,0,0.82}
\definecolor{backcolour}{rgb}{0.95,0.95,0.92}

\lstdefinestyle{mystyle}{
    backgroundcolor=\color{backcolour},   
    commentstyle=\color{codegreen},
    keywordstyle=\color{magenta},
    numberstyle=\tiny\color{codegray},
    stringstyle=\color{codepurple},
    basicstyle=\ttfamily\footnotesize,
    breakatwhitespace=false,         
    breaklines=true,                 
    captionpos=b,                    
    keepspaces=true,                 
    numbers=left,                    
    numbersep=5pt,                  
    showspaces=false,                
    showstringspaces=false,
    showtabs=false,                  
    tabsize=2
}
\lstset{style=mystyle}

\title{\textbf{EV Station Hub}\\ Node-RED OCPP Backend Implementation}
\author{Ruslee Sutthaweekul}
\date{\today}

\begin{document}

\maketitle

\section{Overview}
This document outlines the architecture and implementation steps for creating a Node-RED based backend to serve as an Open Charge Point Protocol (OCPP) server. This server will provide live telemetry data to the Charger-UI and accept remote control commands, specifically the \texttt{Remote Stop} command from the frontend.

\section{Architecture}
The architecture consists of:
\begin{enumerate}
    \item \textbf{Node-RED}: Acts as the Central System (CS) handling OCPP connections and REST API requests.
    \item \textbf{Frontend (Charger-UI)}: Consumes REST endpoints to display data and send commands.
    \item \textbf{Global Store}: Node-RED's global context holds the real-time state of all charging stations.
\end{enumerate}

\vspace{0.5cm}

\tikzstyle{box} = [rectangle, rounded corners, minimum width=3.5cm, minimum height=1.2cm, text centered, text width=3cm, draw=black, thick, fill=gray!10]
\tikzstyle{arrow} = [thick,->,>=stealth]

\begin{figure}[h]
\centering
\begin{tikzpicture}[node distance=3cm]

% Nodes
\node (ui) [box, fill=blue!10] {\textbf{Charger-UI}\\ (Next.js/React)};
\node (nodered) [box, below of=ui, fill=purple!10] {\textbf{Node-RED}\\ (Central System)};
\node (store) [box, right of=nodered, xshift=2.5cm, fill=yellow!10] {\textbf{Global Store}\\ (In-Memory/DB)};
\node (charger) [box, below of=nodered, fill=orange!10] {\textbf{EV Chargers}\\ (Physical Units)};

% Arrows
\draw [arrow] ([xshift=-0.5cm]ui.south) -- node[anchor=east] {\small REST/WS} ([xshift=-0.5cm]nodered.north);
\draw [arrow] ([xshift=0.5cm]nodered.north) -- node[anchor=west] {\small JSON Payload} ([xshift=0.5cm]ui.south);

\draw [arrow] ([yshift=0.2cm]nodered.east) -- node[anchor=south] {\small Write State} ([yshift=0.2cm]store.west);
\draw [arrow] ([yshift=-0.2cm]store.west) -- node[anchor=north] {\small Read State} ([yshift=-0.2cm]nodered.east);

\draw [arrow] ([xshift=-0.5cm]charger.north) -- node[anchor=east] {\small OCPP Telemetry} ([xshift=-0.5cm]nodered.south);
\draw [arrow] ([xshift=0.5cm]nodered.south) -- node[anchor=west] {\small OCPP Commands} ([xshift=0.5cm]charger.north);

\end{tikzpicture}
\caption{System Architecture Flow Diagram}
\end{figure}

\section{Node-RED Flow Configuration}

\subsection{1. Initializing the Station State}
Create an \texttt{Inject} node set to trigger once at startup, connected to a \texttt{Function} node to initialize the state:

\begin{lstlisting}[caption=Initialize Stations]
global.set("ocpp_stations", [
  {
    id: "station-bkk-001",
    name: "Siam Paragon Hub",
    lat: 13.7462,
    lng: 100.5347,
    status: "Available",
    currentMeter: 0,
    powerOutput: 0,
    powerLimit: 120
  },
  {
    id: "station-bkk-002",
    name: "Sukhumvit 21 Fast",
    lat: 13.7410,
    lng: 100.5600,
    status: "Charging",
    currentMeter: 42.15,
    powerOutput: 68.4,
    powerLimit: 150
  }
]);
return msg;
\end{lstlisting}

\subsection{2. REST API: Get Station Data}
To allow the UI to fetch the latest telemetry:
\begin{itemize}
    \item \textbf{HTTP In Node}: Method \texttt{GET}, URL \texttt{/api/ocpp/stations}
    \item \textbf{Function Node}: 
\begin{lstlisting}
msg.payload = global.get("ocpp_stations");
return msg;
\end{lstlisting}
    \item \textbf{HTTP Response Node}: Returns the JSON payload.
\end{itemize}

\subsection{3. REST API: Remote Stop Command}
To accept a "Stop Charging" command from the UI:
\begin{itemize}
    \item \textbf{HTTP In Node}: Method \texttt{POST}, URL \texttt{/api/ocpp/command/stop}
    \item \textbf{Function Node}: Process the incoming request body (e.g., \texttt{\{"stationId": "station-bkk-002"\}})
\begin{lstlisting}[caption=Handle Remote Stop Request]
const stationId = msg.payload.stationId;
let stations = global.get("ocpp_stations") || [];

const targetStation = stations.find(s => s.id === stationId);
if (targetStation && targetStation.status === "Charging") {
    // 1. Update internal simulated state
    targetStation.status = "Available";
    targetStation.powerOutput = 0;
    
    // 2. Real Hardware Integration Point
    // Format OCPP RemoteStopTransaction payload:
    const ocppMsg = {
        payload: {
            command: "RemoteStopTransaction",
            transactionId: targetStation.currentTransactionId // Tracked in state
        }
    };
    // This ocppMsg would be sent to the OCPP out node
    
    global.set("ocpp_stations", stations);
    msg.payload = { success: true, message: "Stop command accepted" };
} else {
    msg.payload = { success: false, error: "Station not charging or not found" };
    msg.statusCode = 400;
}
return msg;
\end{lstlisting}
    \item \textbf{HTTP Response Node}: Returns the success/error response to the UI.
\end{itemize}

\subsection{4. Incoming OCPP: Handling MeterValues}
When a physical charger sends a \texttt{MeterValues} message, Node-RED must intercept this payload and update the global \texttt{ocpp\_stations} store so the UI reflects the live energy, power metrics, and capacity (\texttt{powerLimit} as a numeric integer) immediately.

\begin{itemize}
    \item \textbf{OCPP Server Node}: Listens for incoming \texttt{MeterValues} actions from the active WebSocket connections.
    \item \textbf{Function Node}: Parse the nested array of sampled values to extract energy and power parameters:
\begin{lstlisting}[caption=Process Incoming MeterValues]
// Payload structure depends on node-red-contrib-ocpp formatting
const stationId = msg.chargeBoxIdentity; 
const meterData = msg.payload.meterValue;

let stations = global.get("ocpp_stations") || [];
const targetStation = stations.find(s => s.id === stationId);

if (targetStation && meterData && meterData.length > 0) {
    // Extract the latest sampled values
    const samples = meterData[0].sampledValue;
    
    samples.forEach(sample => {
        if (sample.measurand === "Energy.Active.Import.Register") {
            // Assuming the charger sends data in Wh; converting to kWh
            targetStation.currentMeter = parseFloat(sample.value) / 1000;
        } else if (sample.measurand === "Power.Active.Import") {
            // Assuming the charger sends data in W; converting to kW
            targetStation.powerOutput = parseFloat(sample.value) / 1000;
        }
    });

    // Save the freshly updated state back to the global store
    global.set("ocpp_stations", stations);
}
return msg;
\end{lstlisting}
    \item \textbf{Optional - Broadcast Node}: If you are using WebSockets to pass data to the Next.js UI instead of REST polling, you can link this Function node to a WebSocket Out node to push live updates directly to the map.
\end{itemize}

\section{Connecting Frontend to Node-RED}

In the React/Next.js UI, bind the \textbf{Stop Charging} button to the new endpoint:

\begin{lstlisting}[caption=Frontend Integration (React)]
const handleStopCharging = async (stationId) => {
    try {
        const response = await fetch("http://localhost:1880/api/ocpp/command/stop", {
            method: "POST",
            headers: { "Content-Type": "application/json" },
            body: JSON.stringify({ stationId })
        });
        
        if (response.ok) {
            console.log("Charging stopped successfully");
            // Optionally manually trigger data refresh
        }
    } catch (error) {
        console.error("Failed to send remote command", error);
    }
};
\end{lstlisting}

\section{Real Hardware Integration (Optional)}
To connect real physical chargers:
\begin{enumerate}
    \item Install the \texttt{node-red-contrib-ocpp} node package.
    \item Use the \texttt{OCPP Server} node to accept incoming WebSocket connections.
    \item Route incoming \texttt{MeterValues} and \texttt{StatusNotification} messages to update the \texttt{global.get("ocpp\_stations")} records.
    \item Use the \texttt{OCPP Server} command functionality to emit the \texttt{RemoteStopTransaction} directly to the specified charge point when the REST POST endpoint is triggered.
\end{enumerate}

\end{document}
